\documentclass[PhD]{uclathes}

\usepackage{graphicx}
\usepackage{amsmath}
\usepackage{amssymb}
\usepackage[ruled,vlined,linesnumbered]{algorithm2e}

\title{Thesis Title}
\author{Author Name}
\department{Department Name}
\degreemonth{March}
\degreeyear{2026}

\chair{Committee Chair}
\member{Committee Member 1}
\member{Committee Member 2}
\member{Committee Member 3}

\dedication{Optional dedication text. Remove this command if not needed.}

\acknowledgments{%
Optional acknowledgments text.
}

\abstract{%
Write a concise abstract summarizing the research problem, approach, main results,
and conclusion.
}

\begin{document}

\makeintropages

\chapter{Introduction}

\section{Context and problem}
Describe the broader area, the specific problem addressed in the thesis, and why it matters.

\section{Purpose and expected contributions}
State the thesis goal and summarize the intended contributions.

\section{Research questions}
This thesis addresses the following research questions:

\begin{enumerate}
    \item \textbf{RQ1:} First research question.
    \item \textbf{RQ2:} Second research question.
    \item \textbf{RQ3:} Third research question.
    \item \textbf{RQ4:} Fourth research question.
\end{enumerate}

\section{Scope and limitations}
Clarify what is included in the study, what is excluded, and any important assumptions.

\section{Thesis outline}
Summarize the structure of the remaining chapters.

\chapter{Background}

\section{Foundational concepts}
Introduce the core concepts, definitions, and notation needed for the rest of the thesis.

\section{Technical background}
Explain the main technical ideas, frameworks, or models relevant to the problem.

\section{System or domain background}
Describe the software, hardware, data, or domain-specific systems used in the work.

\chapter{Related Work}

\section{Prior approaches}
Review the most relevant prior literature and group it into clear themes.

\section{Research gap}
Explain what remains unclear or insufficiently explored in prior work.

\chapter{Method}

\section{Research design}
Describe the overall methodology and justify the chosen approach.

\section{Experimental setup}
Document the environment, tools, hardware, datasets, and reproducibility details.

\section{Evaluation criteria}
Define the metrics, comparisons, and criteria used to assess the system or hypothesis.

\section{Model or algorithm specification}
Describe the model, pipeline, or algorithm under study.

\begin{algorithm}[H]
\caption{Example Algorithm Placeholder}
\DontPrintSemicolon
\KwIn{Input data $X$}
\KwOut{Output $Y$}
Initialize the method\;
Process the input\;
Return the result\;
\end{algorithm}

\chapter{System Design}

\section{Overall architecture}
Describe the full system and how the major components interact.

\section{Core components}
Explain the main modules, interfaces, or services.

\section{Implementation decisions}
Document important tradeoffs and design choices.

\chapter{Results}

\section{RQ1 results}
Present findings relevant to the first research question.

\section{RQ2 results}
Present findings relevant to the second research question.

\section{RQ3 results}
Present findings relevant to the third research question.

\section{RQ4 results}
Present findings relevant to the fourth research question.

\chapter{Discussion}

\section{Interpretation of results}
Discuss what the results mean and how they relate to the research questions.

\section{Implications}
Explain the practical or theoretical implications of the findings.

\section{Threats to validity}
Discuss limitations, possible confounders, and how they affect the conclusions.

\chapter{Conclusion}

\section{Summary}
Summarize the main findings and answer the research questions at a high level.

\section{Future work}
Suggest concrete next steps for extending the work.

% Uncomment if needed.
% \appendix
% \chapter{Appendix Title}
% Supplementary material goes here.

% Uncomment if using BibTeX.
% \bibliographystyle{uclathes}
% \bibliography{references}

\end{document}